% !TeX program = latexmk --lualatex --interaction=nonstopmode --halt-on-error --output-directory=build --synctex=1

\documentclass{article}
\usepackage[a4paper, margin=30mm]{geometry}

\usepackage[
    colorlinks,
    citecolor=black,
    linkcolor=black,
]{hyperref}
\usepackage[
  style=numeric,
  doi=true, isbn=false, url=false, eprint=false,
  date=year,% print only the year, not the month
]{biblatex}
\addbibresource{references.bib}

\usepackage{amsmath}
\usepackage{amssymb}
\numberwithin{equation}{section}

\usepackage{cancel}
\usepackage[
    per-mode=fraction
]{siunitx}
    
\usepackage{tikz}

\newcommand{\dd}{\mathrm{d}}
\newcommand{\pd}[2]{\dfrac{\partial #1}{\partial #2}}
\newcommand{\pdd}[2]{\dfrac{\partial^2 #1}{\partial #2^2}}
\newcommand{\dimless}[1]{\hat{#1}}

\setlength{\parindent}{0pt}

\title{Immersed Boundary Conditions}
\author{Gidon Bauer}

\begin{document}

% Name: RWTH Colors
\definecolor{rwth-blue-100}{RGB}{0,84,159}
\definecolor{rwth-blue-75}{RGB}{64,127,183}
\definecolor{rwth-blue-50}{RGB}{142,186,229}
\definecolor{rwth-blue-25}{RGB}{199,221,242}
\definecolor{rwth-blue-10}{RGB}{232,241,250}
\definecolor{rwth-black-100}{RGB}{0,0,0}
\definecolor{rwth-black-75}{RGB}{100,101,103}
\definecolor{rwth-black-50}{RGB}{156,158,159}
\definecolor{rwth-black-25}{RGB}{207,209,210}
\definecolor{rwth-black-10}{RGB}{236,237,237}
\definecolor{rwth-magenta-100}{RGB}{227,0,102}
\definecolor{rwth-magenta-75}{RGB}{233,96,136}
\definecolor{rwth-magenta-50}{RGB}{241,158,177}
\definecolor{rwth-magenta-25}{RGB}{249,210,218}
\definecolor{rwth-magenta-10}{RGB}{253,238,240}
\definecolor{rwth-yellow-100}{RGB}{255,237,0}
\definecolor{rwth-yellow-75}{RGB}{255,240,85}
\definecolor{rwth-yellow-50}{RGB}{255,245,155}
\definecolor{rwth-yellow-25}{RGB}{255,250,209}
\definecolor{rwth-yellow-10}{RGB}{255,253,238}
\definecolor{rwth-petrol-100}{RGB}{0,97,101}
\definecolor{rwth-petrol-75}{RGB}{45,127,131}
\definecolor{rwth-petrol-50}{RGB}{125,164,167}
\definecolor{rwth-petrol-25}{RGB}{191,208,209}
\definecolor{rwth-petrol-10}{RGB}{230,236,236}
\definecolor{rwth-turquoise-100}{RGB}{0,152,161}
\definecolor{rwth-turquoise-75}{RGB}{0,177,183}
\definecolor{rwth-turquoise-50}{RGB}{137,204,207}
\definecolor{rwth-turquoise-25}{RGB}{202,231,231}
\definecolor{rwth-turquoise-10}{RGB}{235,246,246}
\definecolor{rwth-green-100}{RGB}{87,171,39}
\definecolor{rwth-green-75}{RGB}{141,192,96}
\definecolor{rwth-green-50}{RGB}{184,214,152}
\definecolor{rwth-green-25}{RGB}{221,235,206}
\definecolor{rwth-green-10}{RGB}{242,247,236}
\definecolor{rwth-maygreen-100}{RGB}{189,205,0}
\definecolor{rwth-maygreen-75}{RGB}{208,217,92}
\definecolor{rwth-maygreen-50}{RGB}{224,230,154}
\definecolor{rwth-maygreen-25}{RGB}{240,243,208}
\definecolor{rwth-maygreen-10}{RGB}{249,250,237}
\definecolor{rwth-orange-100}{RGB}{246,168,0}
\definecolor{rwth-orange-75}{RGB}{250,190,80}
\definecolor{rwth-orange-50}{RGB}{253,212,143}
\definecolor{rwth-orange-25}{RGB}{254,234,201}
\definecolor{rwth-orange-10}{RGB}{255,247,234}
\definecolor{rwth-red-100}{RGB}{204,7,30}
\definecolor{rwth-red-75}{RGB}{216,92,65}
\definecolor{rwth-red-50}{RGB}{230,150,121}
\definecolor{rwth-red-25}{RGB}{243,205,187}
\definecolor{rwth-red-10}{RGB}{250,235,227}
\definecolor{rwth-bordeaux-100}{RGB}{161,16,53}
\definecolor{rwth-bordeaux-75}{RGB}{182,82,86}
\definecolor{rwth-bordeaux-50}{RGB}{205,139,135}
\definecolor{rwth-bordeaux-25}{RGB}{229,197,192}
\definecolor{rwth-bordeaux-10}{RGB}{245,232,229}
\definecolor{rwth-violet-100}{RGB}{97,33,88}
\definecolor{rwth-violet-75}{RGB}{131,78,117}
\definecolor{rwth-violet-50}{RGB}{168,133,158}
\definecolor{rwth-violet-25}{RGB}{210,192,205}
\definecolor{rwth-violet-10}{RGB}{237,229,234}
\definecolor{rwth-purple-100}{RGB}{122,111,172}
\definecolor{rwth-purple-75}{RGB}{155,145,193}
\definecolor{rwth-purple-50}{RGB}{188,181,215}
\definecolor{rwth-purple-25}{RGB}{222,218,235}
\definecolor{rwth-purple-10}{RGB}{242,240,247}


\maketitle

\section{Immersed Boundary Conditions for Velocity}

Luchini et al. implemented a second order accurate immersed boundary method by correcting the viscous operator in the Navier--Stokes equation to account for the immersed wall \cite{luchiniSimpleEfficientSecondorder2025}.
The Navier--Stokes equations
\begin{equation}
    \begin{gathered}
        \pd{\vec{u}}{t}
        + \left(\vec{u} \cdot \nabla\right) \vec{u}
        + \frac1\rho \nabla p
        - \nu \nabla^2 \vec{u}
        = \vec{f} \\
        \nabla \cdot \vec{u} = 0
    \end{gathered}
    \label{eq:navier-stokes}
\end{equation}
are discretized using second order central finite differences on a MAC-staggered mesh \cite{luchiniSimpleEfficientSecondorder2025,harlowMAC1965}.
Near the wall, it is assumed that the Laplacian from the viscous term becomes dominant \cite{luchiniSimpleEfficientSecondorder2025}.
On a uniform grid, the discretized Laplacian is
\begin{equation}
    \begin{split}
        \nabla^2 u \approx
        \ &d^{(2)}_{x;i,j}(-1) u_{i-1,j}
        + d^{(2)}_{x;i,j}( 0) u_{i,  j}
        + d^{(2)}_{x;i,j}( 1) u_{i+1,j}
        + \\ &d^{(2)}_{y;i,j}(-1) u_{i,j-1}
        + d^{(2)}_{y;i,j}( 0) u_{i,j  }
        + d^{(2)}_{y;i,j}( 1) u_{i,j+1}
    \end{split}
    \label{eq:discrete_laplacian}
\end{equation}

We assume that the immersed wall is to the left of the point $(i,j)$ as shown in Fig. \ref{fig:ib-to-left}.
The value at point $(i-1,j)$ is extrapolated linearly such that the value at the interface is exactly $u^*$.
\begin{equation}
    u_{i-1,j} = u^* + \frac{u_{i,j} - u^*}{\delta x} (\delta x - \Delta x)
    = u_{i,j} - \frac{(u_{i,j} - u^*) \Delta x}{\delta x}
    = u_{i,j} \left(1 - \frac{\Delta x}{\delta x}\right) + u^* \frac{\Delta x}{\delta x}
\end{equation}
Assuming a stationary wall, we simply get
\begin{equation}
    u_{i-1,j} = u_{i,j} \left(1 - \frac{\Delta x}{\delta x}\right).
\end{equation}

We then insert this into Eq. \eqref{eq:discrete_laplacian} and get
\begin{equation}
    \begin{split}
        \nabla^2 u \approx
        \ &d^{(2)}_{x;i,j}(-1) u_{i,j} \left(1 - \frac{\Delta x}{\delta x}\right)
        + d^{(2)}_{x;i,j}( 0) u_{i,  j}
        + d^{(2)}_{x;i,j}( 1) u_{i+1,j}
        + \\ &d^{(2)}_{y;i,j}(-1) u_{i,j-1}
        + d^{(2)}_{y;i,j}( 0) u_{i,j  }
        + d^{(2)}_{y;i,j}( 1) u_{i,j+1}
    \end{split}
\end{equation}
This change can be absorbed into the factor $d^{(2)}_{x;i,j}(0)$ in the following way
\begin{equation}
    \tilde{d}^{(2)}_{x;i,j}(0) =
    d^{(2)}_{x;i,j}(0) -
    \underbrace{d^{(2)}_{x;i,j}(-1) \left(\frac{\Delta x}{\delta x} - 1\right)}_{=: \lambda}.
\end{equation}
We can define the relative distance $\xi = \frac{\delta x}{\Delta x} \in (0, 1]$ and get
\begin{equation}
    \lambda = d^{(2)}_{x;i,j}(-1) \frac{1 - \xi}{\xi}.
    \label{eq:calc-lambda}
\end{equation}

\begin{figure}
    \centering
    \begin{tikzpicture}[>=stealth, scale=2, line width=2]
        \draw (0, 0) -- (3, 0);
        \draw (0, 1) -- (3, 1);
        \draw (0, 2) -- (3, 2);
        \draw (0, 3) -- (3, 3);
        \draw (0, 0) -- (0, 3);
        \draw (1, 0) -- (1, 3);
        \draw (2, 0) -- (2, 3);
        \draw (3, 0) -- (3, 3);
        \node[below left] at (0, 1.5) {$(i-1,j)$};
        \node[below right] at (1, 1.5) {$(i,j)$};
        
        \draw (0.2, -0.1) -- (0.6, 3.1);

        \draw[->, rwth-red-100] (-0.2, 1.5) -- (0.2, 1.5);
        \draw[->, rwth-red-100] ( 0.8, 1.5) -- (1.2, 1.5);
        \draw[->, rwth-red-100] ( 1.8, 1.5) -- (2.2, 1.5);
        \draw[->, rwth-red-100] ( 0.8, 0.5) -- (1.2, 0.5);
        \draw[->, rwth-red-100] ( 0.8, 2.5) -- (1.2, 2.5);

        \draw[|-|, line width=1.5, rwth-purple-100] (0, 1.5) -- (0.4, 1.5);
        \node[below, rwth-purple-100] at (0.2, 1.5) {$\overline{\delta x}$};
        \draw[|-|, line width=1.5, rwth-purple-100] (0.4, 1.5) -- (1, 1.5);
        \node[below, rwth-purple-100] at (0.7, 1.5) {$\delta x$};
    \end{tikzpicture}
    \caption{Immersed boundary to the left of point $(i,j)$. Here $\overline{\delta x} = \Delta x - \delta x$.}
    \label{fig:ib-to-left}
\end{figure}

\subsection{General Orthogonal Coordinate System}

In a general three-dimensional orthogonal coordinate system $(q_1, q_2, q_3)$ with corresponding scale factors $h_1, h_2, h_3$,
the scalar Laplacian is
\begin{equation}
    \nabla^2 u = \frac{1}{H} \left[
        \pd{}{q_1} \left(\frac{H}{h_1^2} \pd{u}{q_1}\right)
        + \pd{}{q_2} \left(\frac{H}{h_2^2} \pd{u}{q_2}\right)
        + \pd{}{q_3} \left(\frac{H}{h_3^2} \pd{u}{q_3}\right)
    \right].
\end{equation}
Here $H = h_1 h_2 h_3$ is the volume metric.

Considering only the $q_1$ direction, we discretize on a uniform mesh
\begin{equation}
    \begin{split}
        \frac{1}{H} \pd{}{q_1} \left(\frac{H}{h_1^2} \pd{u}{q_1}\right)\bigg\vert_{i} &\approx
        \frac{1}{H_{i}} \frac{1}{\Delta q_1} \left(
            \frac{H_{i+\tfrac12}}{h_{1, i+\tfrac12}^2} \frac{u_{i+1} - u_{i}}{\Delta q_1}
            - \frac{H_{i-\tfrac12}}{h_{1, i-\tfrac12}^2} \frac{u_{i} - u_{i-1}}{\Delta q_1}
        \right) + K \\
        &= \frac{1}{H_i \Delta q_1^2} \left[
            \frac{H_{i-\tfrac12}}{h_{1,i-\tfrac12}^2} u_{i-1} -
            \left(
                \frac{H_{i+\tfrac12}}{h_{1,i+\tfrac12}^2} +
                \frac{H_{i-\tfrac12}}{h_{1,i-\tfrac12}^2}
            \right) u_{i} +
            \frac{H_{i+\tfrac12}}{h_{1,i+\tfrac12}^2} u_{i+1}
        \right] + K
    \end{split}
\end{equation}
Here $K$ describes the coupling terms that arise for the vector Laplacian in orthogonal coordinates.
These terms are not relevant for the subsequent derivations.

This gives us the values for $d^{(2)}_{q_k;i,j}$:
\begin{equation}
    \begin{split}
        d^{(2)}_{q_k;i,j}(-1) &= \frac{1}{H_i \Delta q_k^2} \frac{H_{i-\tfrac12}}{h_{k,i-\tfrac12}^2} \\
        d^{(2)}_{q_k;i,j}(0)  &= -\frac{1}{H_i \Delta q_k^2} \left(
                \frac{H_{i+\tfrac12}}{h_{k,i+\tfrac12}^2} +
                \frac{H_{i-\tfrac12}}{h_{k,i-\tfrac12}^2}
            \right) \\
        d^{(2)}_{q_k;i,j}(1)  &= \frac{1}{H_i \Delta q_k^2} \frac{H_{i+\tfrac12}}{h_{k,i+\tfrac12}^2}
    \end{split}
\end{equation}
Note that for $k=2$, the index must be changed to $j$.

\subsubsection{Cartesian Coordinates}

For Cartesian coordinates $(q_1,q_2,q_3) = (x,y,z)$, we have $h_1 = h_2 = h_3 = 1$.
Then we have for the parameters of the discretized Laplacian
\begin{equation}
    \begin{split}
        d^{(2)}_{q_k;i,j}(-1) &=  \frac{1}{\Delta q_k^2} \\
        d^{(2)}_{q_k;i,j}(0)  &= -\frac{2}{\Delta q_k^2} \\
        d^{(2)}_{q_k;i,j}(1)  &=  \frac{1}{\Delta q_k^2}.
    \end{split}
\end{equation}

\subsubsection{Polar Coordinates}

For polar coordinates in two dimensions $(q_1,q_2) = (\theta,r)$, we have $h_1 = r$ and $h_2 = 1$.
Then we have for the parameters of the discretized Laplacian
\begin{equation}
    \begin{split}
        d^{(2)}_{\theta;i,j}(-1) &= \frac{1}{r_i \Delta \theta^2} \frac{1}{r_{i-\tfrac12}} \\
        d^{(2)}_{\theta;i,j}(0)  &= -\frac{1}{r_i \Delta \theta^2} \left(
                \frac{1}{r_{i+\tfrac12}} +
                \frac{1}{r_{i-\tfrac12}}
            \right) \\
        d^{(2)}_{\theta;i,j}(1)  &= \frac{1}{r_i \Delta \theta^2} \frac{1}{r_{i+\tfrac12}} \\
        d^{(2)}_{r;i,j}(-1) &= \frac{1}{r_i \Delta r^2} r_{i-\tfrac12} \\
        d^{(2)}_{r;i,j}(0)  &= -\frac{1}{r_i \Delta r^2} \left(r_{i+\tfrac12} + r_{i-\tfrac12}\right) \\
        d^{(2)}_{r;i,j}(1)  &= \frac{1}{r_i \Delta r^2} r_{i+\tfrac12}
    \end{split}
\end{equation}

\subsection{Time Integration}

The spatial derivatives of Eq. \eqref{eq:navier-stokes} are discretized with the correction for the immersed boundary as described above.
\begin{equation}
    \pd{u_{i,j}}{t}
    = -\lambda_{i,j} \nu u_{i,j} +  R_{i,j}
    = -\lambda_{i,j} \nu u_{i,j}
    + d^{(2)}_{x;i,j}( 0) u_{i,  j}
    + d^{(2)}_{x;i,j}( 1) u_{i+1,j}
    + \dots
    \label{eq:ODE-update}
\end{equation}
The variable $R_{i,j}$ combines all spatial terms, only the first part of the discretized Laplacian is shown explicitly \cite{luchiniSimpleEfficientSecondorder2025}.
We again assume that the immersed interface is located to the left of point $(i,j)$.

This becomes a stiff problem whenever the immersed interface is close to a grid point.
Then $\delta x \to 0$ and $\lambda \to \infty$.
Therefore, the regular, uncorrected right-hand-side is treated explicitly but the correction implicitly.

Using the implicit Euler discretization for the correction, we obtain
\begin{equation}
    \begin{split}
        &\frac{u_{i,j}^{n+1} - u_{i,j}^n}{\Delta t} =
        -\lambda_{i,j} \nu u_{i,j}^{n+1} + R_{i,j}^n \\
        \Rightarrow\ 
        &u_{i,j}^{n+1} =
        \frac{u_{i,j}^n + \Delta t R_{i,j}^n}{1 + \Delta t \lambda_{i,j} \nu}
    \end{split}
\end{equation}
which is a simple addition to the vanilla numerical method.

For higher-accuracy, Luchini et al. use the analytical solution of the ODE \eqref{eq:ODE-update}.
They assume that $R_{i,j}$ is constant during the integration.
They then get the general update rule
\begin{equation}
    \begin{split}
        &A u_{i,j}^{n+1} - B u_{i,j}^n = C F_{i,j}^n \\
        \Rightarrow\ 
        &u_{i,j}^{n+1} = \frac{B u_{i,j}^n + C F_{i,j}^n}{A}
    \end{split}
    \label{eq:analytical-correction}
\end{equation}
with
\begin{equation}
    \begin{split}
        &B = \frac{\lambda_{i,j} \nu \Delta t}{\exp(\lambda_{i,j} \nu \Delta t) - 1} \\
        &A = \lambda_{i,j} \nu \Delta t + B \\
        &C = \Delta t.
    \end{split}
\end{equation}
The variable $F_{i,j}^n$ is a placeholder for the explicit evaluation of the right-hand-side at any time level as the method requires.
In the simplest case we have $F_{i,j}^n = R_{i,j}^n$.

\subsection{Further Remarks}

\begin{itemize}
    \item The velocity values inside of the solid must be set to zero.
    \item The corrections are linear and can be calculated for each direction separately and then added.
    \item The calculated value for the velocity is the predicted value in the prediction-correction scheme.
    This has to be corrected for it to be divergence-free.
    The pressure Poisson equation is not changed.
\end{itemize}

\subsection{Moving Walls}

For moving walls we have $u^* \ne 0$.
The discretized Laplacian then becomes
\begin{equation}
    \begin{split}
        \nabla^2 u \approx
        \ &d^{(2)}_{x;i,j}(-1) \left[
            u_{i,j} \left(1 - \frac{\Delta x}{\delta x}\right) +
            u^* \frac{\Delta x}{\delta x}
        \right]
        + d^{(2)}_{x;i,j}( 0) u_{i,  j}
        + d^{(2)}_{x;i,j}( 1) u_{i+1,j}
        + \dots
    \end{split}
\end{equation}

The term $d^{(2)}_{x;i,j}(-1) u^* \frac{\Delta x}{\delta x}$ cannot be absorbed into $d^{(2)}_{x;i,j}(0)$.
We therefore have an additional term in the correction $\eta = d^{(2)}_{x;i,j}(-1) \frac{u^*}{\xi}$.

\begin{equation}
    \begin{split}
        \pd{u_{i,j}}{t}
        &= -\lambda_{i,j} \nu u_{i,j} + \nu \eta_{i,j} +  R_{i,j} \\
        &= -\lambda_{i,j} \nu u_{i,j} + \nu \eta_{i,j}
        + d^{(2)}_{x;i,j}( 0) u_{i,  j}
        + d^{(2)}_{x;i,j}( 1) u_{i+1,j}
        + \dots
    \end{split}
\end{equation}

\textcolor{rwth-red-100}{\texttt{[TODO]} This is not complete yet.}

% ==================================================================================================
\section{Immersed Boundary Conditions for Species Transport}

The conservation of a passive scalar $s$ in an incompressible flow is written as
\begin{equation}
    \pd{s}{t} + (\vec{u} \cdot \nabla) s - D \nabla^2 s = 0.
    \label{eq:species}
\end{equation}
Here, $D$ is the diffusion coefficient which is assumed to be constant.
We can apply the same correction as for the velocity for the scalar transport.
Now we need to use the diffusion coefficient $D$ instead of the kinematic viscosity $\nu$ and
the correction must be calculated for the centered cells.

\subsection{Homogeneous Dirichlet Boundary Conditions}

Applying the same derivation to Eq. \eqref{eq:species}, we get the Dirichlet boundary conditions $s = 0$.
We still calculate $\lambda$ via Eq. \eqref{eq:calc-lambda}.

In the implicit Euler time discretization, we have
\begin{equation}
    s_{i,j}^{n+1} =
    \frac{s_{i,j}^n + \Delta t R_{i,j}^n}{1 + \Delta t D \lambda_{i,j}}
\end{equation}
where $R$ are now the discretized spatial derivatives of Eq. \eqref{eq:species} moved to the right-hand-side.

The analytical correction Eq. \eqref{eq:analytical-correction} has parameters
\begin{equation}
    \begin{split}
        &B = \frac{\lambda_{i,j} D \Delta t}{\exp(\lambda_{i,j} D \Delta t) - 1} \\
        &A = \lambda_{i,j} D \Delta t + B \\
        &C = \Delta t
    \end{split}
\end{equation}

\subsection{Neumann Boundary Conditions}

For most applications, homogeneous Dirichlet boundary conditions are not the right choice for the species equation.
One usually chooses Neumann boundary conditions
\begin{equation}
    \nabla s \cdot \vec{n} = v.
\end{equation}
Here $\vec{n}$ is the outwards pointing normal and $v$ a chosen value.

These kind of boundaries are not easily implemented with the method by Luchini et al.

\begin{figure}[h]
    \centering
    \begin{tikzpicture}[>=stealth, scale=2, line width=2]
        \draw (0, 0) -- (3, 0);
        \draw (0, 1) -- (3, 1);
        \draw (0, 2) -- (3, 2);
        \draw (0, 3) -- (3, 3);
        \draw (0, 0) -- (0, 3);
        \draw (1, 0) -- (1, 3);
        \draw (2, 0) -- (2, 3);
        \draw (3, 0) -- (3, 3);
        \node[below left] at (0, 1.5) {$(i-1,j)$};
        \node[below right] at (1, 1.5) {$(i,j)$};
        
        \draw[rwth-black-75] (0.2, -0.1) -- (0.6, 3.1);
        \draw[->, rwth-black-75] (0.4, 1.5) -- (0.4-3*0.99/4, 1.5+3*0.12/4) node [above] {$\vec{n}$};

        \draw[->, rwth-red-100] (-0.2, 1.5) -- (0.2, 1.5);
        \draw[->, rwth-red-100] ( 0.8, 1.5) -- (1.2, 1.5);
        \draw[->, rwth-red-100] ( 1.8, 1.5) -- (2.2, 1.5);
        \draw[->, rwth-red-100] ( 0.8, 0.5) -- (1.2, 0.5);
        \draw[->, rwth-red-100] ( 0.8, 2.5) -- (1.2, 2.5);

        \draw[|-|, line width=1.5, rwth-purple-100] (0, 1.5) -- (0.4, 1.5);
        \node[below, rwth-purple-100] at (0.2, 1.5) {$\overline{\delta x}$};
        \draw[|-|, line width=1.5, rwth-purple-100] (0.4, 1.5) -- (1, 1.5);
        \node[below, rwth-purple-100] at (0.7, 1.5) {$\delta x$};
    \end{tikzpicture}
    \caption{
        Immersed boundary to the left of point $(i,j)$.
        Here $\overline{\delta x} = \Delta x - \delta x$.
        The vector $\vec{n}$ points into the solid.
    }
    \label{fig:ib-to-left-neumann}
\end{figure}

\newpage
\printbibliography

\end{document}